\documentclass{article}

\textwidth=6.5in
\textheight=8.5in
\voffset=-54pt
\oddsidemargin=0in
\evensidemargin=0in
\setlength{\parskip}{8pt}
\setlength{\parindent}{0pt}

\usepackage{workbook}

\usepackage[sols]{handout}

\begin{document}

\begin{center}
  \large\textsc{Problem Set 6 - Circular Motion, Arc Length, and Angles}
  
      \pspace{12pt}
  \begin{problemonly}\large Name: \rule{5cm}{0.15mm}\\ \end{problemonly}
\end{center}


\begin{enumerate}

\item 
\begin{problem}
    In an amusement park, a Ferris wheel with a radius of 180 feet makes one revolution every 3 minutes.
\end{problem}

\begin{enumerate}
    \item 
    \begin{problem}
    Determine the (linear) speed of the riders on the Ferris wheel in feet per second. 
    \end{problem}
    \pspace{\vfill}

    \begin{solution}
        The distance that a rider travels in one revolution is $2\pi(180\text{ ft})$ or $360\pi$ feet. Since the Ferris wheel makes one revolution every 3 minutes, the rider's speed is $\frac{360\pi\text{ ft}}{180\text{ s}}$ or \fbox{$2\pi$ ft/s.}
    \end{solution}

    \item 
    \begin{problem}
        At what rate does the Ferris wheel rotate (called its \emph{angular speed}) in degrees per second?
    \end{problem}
    \pspace{\vfill}

    \begin{solution}
        In one revolution, the Ferris wheel rotates \ang{360}. The Ferris wheel's angular speed is therefore $\frac{\ang{360}}{180\text{ s}}$ or \fbox{$2$ degrees per second.}
    \end{solution}



    \item 
    \begin{problem}
        What is the Ferris wheel's angular speed in radians per second?
    \end{problem}
    \pspace{\vfill}

    \begin{solution}
        In one revolution, the Ferris wheel rotates $2\pi$ radians. The Ferris wheel's angular speed is therefore $\frac{2\pi\text{ rad}}{180\text{ s}}$ or \fbox{$\frac{\pi}{180}$ radians per second.}
    \end{solution}

    \item
    \begin{problem}
        When a rider on the Ferris wheel has traveled a distance of 180 feet, what angle has the Ferris wheel rotated by? (Hint: one angle unit has a special connection to radius and arc length.)
    \end{problem}
    \pspace{\vfill}

    \begin{solution}
The relationship between arc length, radius, and angle is $s=r\theta$, where $s$ is the arc length, $r$ is the radius, and $\theta$ is the angle in radians. Since the radius of the Ferris wheel is 180 feet, if the arc length traveled by the rider is also 180 feet, this means that the Ferris wheel has rotated by an angle of \fbox{1 radian.}
    \end{solution}

    \item 
    \begin{problem}
        The amusement park wants to build a larger Ferris wheel that still completes one revolution every 3 minutes. If the new Ferris wheel is twice as large in diameter, what will happen to the i. linear speed and ii. angular speed experienced by the riders? 
    \end{problem}
    \pspace{\vfill}

    \begin{solution}
        If the new Ferris wheel is twice as large in diameter, then the circumference will also be twice as large. If the Ferris wheel still completes one revolution every 3 minutes, this means that the riders will travel twice the distance in the same amount of time, so the linear speed will be twice as large.

        No matter the radius, one revolution is always an angle of \ang{360} or $2\pi\text{ rad}$. This means that the angular speed will not change.
    \end{solution}
\end{enumerate}

\pspace{\newpage}

\item 
\begin{problem}
    You can make a conical cup by cutting a \emph{sector} (a wedge) out of a circular piece of paper and joining the edges together, as shown below. Suppose the radius of the circular paper is 5 cm and $\theta = 6\pi/5$.
\end{problem}
\begin{center}
    {\setlength{\tabcolsep}{0.5in}
    \begin{tabular}{cc}
\begin{tikzpicture}[scale=0.3]

  % Circle center
  \coordinate (O) at (0,0);

  % Draw the circle with radius 6 cm
  \draw[very thick] (O) -- ++(0:6) arc[start angle=0,end angle=-216,radius=6] -- cycle;

  \draw (1.5,0) arc[start angle=0,end angle=-216,radius=1.5];

  % % Draw radius lines
  % \draw[thick] (O) -- (6,0); % Initial radius
  % \draw[thick] (O) -- (-3,5.196); % Terminal radius, calculated from the angle
  \node at (270:2.3) {$\theta$};

  % Label the radius
  \node at (3, 0.7) {5 cm};
  \node at (-1, 2.5) {5 cm};

\end{tikzpicture} &
  \begin{tikzpicture}[scale=0.6]
    \draw[thick] (0,0) arc (170:10:2cm and 0.8cm)coordinate[pos=0] (a);
    \draw[thick] (0,0) arc (-170:-10:2cm and 0.8cm)coordinate (b);
    \draw[densely dashed] ([yshift=-4cm]$(a)!0.5!(b)$) -- node[right,font=\footnotesize] {$h$}coordinate[pos=0.95] (aa)($(a)!0.5!(b)$)
                            -- node[above,font=\footnotesize] {$r$}coordinate[pos=0.1] (bb) (b);
    \draw[thick] (aa) -| (bb);
    \draw[thick] (a) -- ([yshift=-4cm]$(a)!0.5!(b)$) -- node[midway,right] {5 cm} (b);
  \end{tikzpicture}
    \end{tabular}
    }
    \end{center}

\begin{enumerate}
    \item 
    \begin{problem}
        Find the circumference of the opening of the cup.
    \end{problem}
    \pspace{\vfill}

    \begin{solution}
        The circumference of the opening of the cup is the arc length of a circular arc with a central angle of $\theta = 6\pi/5$ and a radius of 5 cm. Using the formula $s=r\theta$, the length is $s=(5\text{ cm})(\frac{6\pi}{5})=\fbox{6\pi\text{ cm}.}$
    \end{solution}

    \item 
    \begin{problem}
        Find the radius $r$ of the opening of the cup. (Hint: use $C=2\pi r$.)
    \end{problem}
    \pspace{\vfill}

    \begin{solution}
        Since the opening of the cup makes a circle with a circumference of $6\pi$ cm, the radius of the opening must be \fbox{$3$ cm.}
    \end{solution}

    \item 
    \begin{problem}
        Find the height $h$ of the cup. (Hint: Use the Pythagorean Theorem.)
    \end{problem}
    \pspace{\vfill}

    \begin{solution}
        Inside the cone, we have a right triangle whose hypotenuse is along the slanted side of the cone, which has a length of 5 cm.
        One of the legs of the right triangle is the radius of the opening, which we determined has a length of $3$ cm.
        The other leg goes from the tip of the cone to the center of the opening, whose length $h$ is the height of the cup.
        \begin{align*}
            h^2 + (3\text{ cm})^2 &= (5\text{ cm})^2 \\
            h &= \sqrt{16 \text{ cm}^2} \\
            h &= \fbox{4\text{ cm}}
        \end{align*}
    \end{solution}

    \item 
    \begin{problem}
        Find the volume of the cup. (Reminder: $V_\text{cone} = \tfrac{1}{3}\pi r^2 h$.)
    \end{problem}
    \pspace{\vfill}

    \begin{solution}
    Combining the information we found, $r=3\text{ cm}$ and $h=4\text{ cm}$:
    \begin{align*}
        V&=\frac{1}{3}\pi(3\text{ cm})^2(4\text{ cm})\\
        &=\fbox{$12\pi \text{ cm}^3$}
        \end{align*}
    \end{solution}
\end{enumerate}

\item 
\begin{problem}
    How many degrees is $6\pi/5$ radians?
\end{problem}
\pspace{\vfill}

\begin{solution}
    To convert $6\pi/5\text{ rad}$ from radians to degrees, we can use a conversion factor. $180$ degrees is the same as $\pi$ radians, so multiplying by $\dfrac{180\text{ deg}}{\pi\text{ rad}}$ is the same as multiplying by 1.
    \[\left(\frac{6\pi}{5}\text{ rad}\right) \times \left(\frac{180\text{ deg}}{\pi\text{ rad}}\right)=\fbox{216\text{ deg}}\]
\end{solution}

\pspace{\newpage}

\item
\begin{problem}
    Find the missing angle measures in each column. One column has been done for you as an example. (r.a. is short for right angle.)
\end{problem}

{\setlength{\tabcolsep}{9pt}
\begin{tabular}{l|l|l|l|l|l}

\cfitb[$\vphantom{\vspace{16pt}}$]{0.5in}{$4$} r.a. &
\cfitb[$\vphantom{\vspace{16pt}}$]{0.5in}{$2$} r.a. &
\cfitb[$\vphantom{\vspace{16pt}}1$]{0.5in}{$\vphantom{\vspace{16pt}}1$} r.a. &
\cfitb[$\vphantom{\vspace{16pt}}$]{0.5in}{$2/3$} r.a. &
\cfitb[$1/2$]{0.5in}{$1/2$} r.a. &
\cfitb[$\vphantom{\vspace{16pt}}$]{0.5in}{$\vphantom{\vspace{16pt}}1/3$} r.a. \\[12pt]

\cfitb[$\vphantom{\vspace{16pt}}360$]{0.5in}{$360$} deg &
\cfitb[$\vphantom{\vspace{16pt}}$]{0.5in}{$180$} deg &
\cfitb[$\vphantom{\vspace{16pt}}90$]{0.5in}{$90$} deg &
\cfitb[$\vphantom{\vspace{16pt}}60$]{0.5in}{$60$} deg &
\cfitb[$\vphantom{\vspace{16pt}}$]{0.5in}{$45$} deg &
\cfitb[$\vphantom{\vspace{16pt}}$]{0.5in}{$30$} deg\\[12pt]

\cfitb[$\vphantom{\vspace{16pt}}$]{0.5in}{$2\pi$} rad & 
\cfitb[$\vphantom{\vspace{16pt}}\pi$]{0.5in}{$\pi$} rad & 
\cfitb[$\vphantom{\vspace{16pt}}\pi/2$]{0.5in}{$\pi/2$} rad &
\cfitb[$\vphantom{\vspace{16pt}}$]{0.5in}{$\pi/3$} rad & 
\cfitb[$\vphantom{\vspace{16pt}}$]{0.5in}{$\pi/4$} rad & 
\cfitb[$\pi/6$]{0.5in}{$\pi/6$} rad

\end{tabular}
}
\pspace{\vfill}

\item 
\begin{problem}
    The \emph{quadrantal angles} are the multiples of $\frac{\pi}{2}$, like $-\frac{\pi}{2}$, $0$, $\frac{\pi}{2}$, $\pi$, $\frac{3\pi}{2}$, $2\pi$, and $\frac{5\pi}{2}$. They divide the $xy$-plane into four quadrants that are conventionally numbered I, II, III, and IV as shown.
    \end{problem}

    \begin{center}
    \begin{tikzpicture}[scale=0.5]
    % Draw the coordinate axes
    \draw[->] (-2,0) -- (2,0) node[right] {$x$};
    \draw[->] (0,-2) -- (0,2) node[above] {$y$};

    % Label quadrants
    \node at (1,1) {\Large I};
    \node at (-1,1) {\Large II};
    \node at (-1,-1) {\Large III};
    \node at (1,-1) {\Large IV};
\end{tikzpicture}
    \end{center}

\begin{problem}
    Shown below are six angles drawn in standard position. For each, determine the quadrant the terminal ray lands in and find the neighboring quadrantal angles. Part (a) has been done as an example.
\end{problem}

\begin{center}
\begin{tabular}{ccc}
(a)\hspace{1em} & (b)\hspace{1em} & (c)\hspace{1em} \\
\begin{tikzpicture}
    \begin{axis}[
        axis equal image,
        xmin=-5, xmax=5, ymin=-5, ymax=5,
        xlabel={$x$},
        ylabel={$y$},
        xtick=\empty, ytick=\empty,
        width=2.5in
      ]
      
      % Draw the angle
      \draw[thick,->] (0,0) -- (\rtheta{5}{0});
      \draw[thick,->] (0,0) -- (\rtheta{5}{155});

      % Draw angle arc
      \addplot[domain=0:155,-stealth, very thick] ({cos(x)},{sin(x)});

      % Label the angle
      \node at (\rtheta{1.75}{155/2}) {$A$};

    \end{axis}
\end{tikzpicture} &
\begin{tikzpicture}
    \begin{axis}[
        axis equal image,
        xmin=-5, xmax=5, ymin=-5, ymax=5,
        xlabel={$x$},
        ylabel={$y$},
        xtick=\empty, ytick=\empty,
        width=2.5in
      ]
      
      % Draw the angle
      \draw[thick,->] (0,0) -- (\rtheta{5}{0});
      \draw[thick,->] (0,0) -- (\rtheta{5}{75});

      % Draw angle arc
      \addplot[domain=0:75,-stealth, very thick] ({cos(x)},{sin(x)});

      % Label the angle
      \node at (\rtheta{1.75}{75/2}) {$B$};

    \end{axis}
\end{tikzpicture}&
\begin{tikzpicture}
    \begin{axis}[
        axis equal image,
        xmin=-5, xmax=5, ymin=-5, ymax=5,
        xlabel={$x$},
        ylabel={$y$},
        xtick=\empty, ytick=\empty,
        width=2.5in
      ]
      
      % Draw the angle
      \draw[thick,->] (0,0) -- (\rtheta{5}{0});
      \draw[thick,->] (0,0) -- (\rtheta{5}{320});

      % Draw angle arc
      \addplot[domain=0:320,-stealth, very thick] ({cos(x)},{sin(x)});

      % Label the angle
      \node at (\rtheta{1.75}{320/2}) {$C$};

    \end{axis}
\end{tikzpicture}\\

$\cfitb[\vphantom{\vspace{16pt}}\frac{\pi}{2}]{0.5in}{\frac{\pi}{2}} < A < \cfitb[\pi\vphantom{\vspace{16pt}}]{0.5in}{\pi}$ \hspace{1em} &
$\cfitb[\vphantom{\vspace{16pt}}]{0.5in}{0} < B < \cfitb[\vphantom{\vspace{16pt}}]{0.5in}{\frac{\pi}{2}}$ \hspace{1em} &
$\cfitb[\vphantom{\vspace{16pt}}]{0.5in}{\frac{3\pi}{2}} < C < \cfitb[\vphantom{\vspace{16pt}}]{0.5in}{2\pi}$ \hspace{1em} \\[12pt]

Quadrant \cfitb[II]{0.5in}{II} & Quadrant \cfitb{0.5in}{I} & Quadrant \cfitb{0.5in}{III}

\end{tabular}

\vspace{16pt}

\begin{tabular}{ccc}
(d)\hspace{1em} & (e)\hspace{1em} & (f)\hspace{1em} \\
\begin{tikzpicture}
    \begin{axis}[
        axis equal image,
        xmin=-5, xmax=5, ymin=-5, ymax=5,
        xlabel={$x$},
        ylabel={$y$},
        xtick=\empty, ytick=\empty,
        width=2.5in
      ]
      
      % Draw the angle
      \draw[thick,->] (0,0) -- (\rtheta{5}{0});
      \draw[thick,->] (0,0) -- (\rtheta{5}{-80});

      % Draw angle arc
      \addplot[domain=-80:0,stealth-, very thick] ({cos(x)},{sin(x)});

      % Label the angle
      \node at (\rtheta{1.75}{-80/2}) {$D$};

    \end{axis}
\end{tikzpicture} &
\begin{tikzpicture}
    \begin{axis}[
        axis equal image,
        xmin=-5, xmax=5, ymin=-5, ymax=5,
        xlabel={$x$},
        ylabel={$y$},
        xtick=\empty, ytick=\empty,
        width=2.5in
      ]
      
      % Draw the angle
      \draw[thick,->] (0,0) -- (\rtheta{5}{0});
      \draw[thick,->] (0,0) -- (\rtheta{5}{395});

      % Draw angle arc
      \addplot[domain=0:395,-stealth, very thick] ({cos(x)*(0.5+x/395)},{sin(x)*(0.5+x/395)});

      % Label the angle
      \node at (\rtheta{1.75}{420/2}) {$E$};

    \end{axis}
\end{tikzpicture} &
\begin{tikzpicture}
    \begin{axis}[
        axis equal image,
        xmin=-5, xmax=5, ymin=-5, ymax=5,
        xlabel={$x$},
        ylabel={$y$},
        xtick=\empty, ytick=\empty,
        width=2.5in
      ]
      
      % Draw the angle
      \draw[thick,->] (0,0) -- (\rtheta{5}{0});
      \draw[thick,->] (0,0) -- (\rtheta{5}{-220});

      % Draw angle arc
      \addplot[domain=-220:0,stealth-, very thick] ({cos(x)},{sin(x)});

      % Label the angle
      \node at (\rtheta{1.75}{-220/2}) {$F$};

    \end{axis}
\end{tikzpicture}\\

$\cfitb[\vphantom{\vspace{16pt}}]{0.5in}{-\frac{\pi}{2}} < D < \cfitb[\vphantom{\vspace{16pt}}]{0.5in}{0}$ \hspace{1em}&
$\cfitb[\vphantom{\vspace{16pt}}]{0.5in}{2\pi} < E < \cfitb[\vphantom{\vspace{16pt}}]{0.5in}{\frac{5\pi}{2}}$ \hspace{1em} &
$\cfitb[\vphantom{\vspace{16pt}}]{0.5in}{\frac{-3\pi}{2}} < F < \cfitb[\vphantom{\vspace{16pt}}]{0.5in}{-\pi}$ \hspace{1em} \\[12pt]

Quadrant \cfitb{0.5in}{IV} & Quadrant \cfitb{0.5in}{I} & Quadrant \cfitb{0.5in}{II}

\end{tabular}
\end{center}

\pspace{\newpage}

\item 
\begin{problem}
    Recall from class that the \emph{reference angle} for an angle $\theta$ is an acute angle that measures how far $\theta$ is from a multiple of $\pi$. Find the reference angle for each angle below, and use it to sketch the angle in standard position.
\end{problem}

\begin{center}
\begin{tabular}{ccc}
(a)\hspace{1em} & (b)\hspace{1em} & (c)\hspace{1em} \\
\begin{tikzpicture}
    \begin{axis}[
        axis equal image,
        xmin=-5, xmax=5, ymin=-5, ymax=5,
        xlabel={$x$},
        ylabel={$y$},
        xtick=\empty, ytick=\empty,
        width=2.5in
      ]
      % Draw the angle
      \draw[thick,->] (0,0) -- (\rtheta{5}{0});
      \draw[thick,->] (0,0) -- (\rtheta{5}{150});

      % Draw angle arc
      \addplot[domain=0:150,-stealth, very thick] ({cos(x)},{sin(x)});

      % Label the angle
      \node at (\rtheta{1.5}{150/2}) {$\theta$};

      % Label the reference angle
      \node at (\rtheta{2}{170}) {\footnotesize$\pi/6$};

    \end{axis}
\end{tikzpicture} &
\begin{tikzpicture}
    \begin{axis}[
        axis equal image,
        xmin=-5, xmax=5, ymin=-5, ymax=5,
        xlabel={$x$},
        ylabel={$y$},
        xtick=\empty, ytick=\empty,
        width=2.5in
      ]
      \addplot[draw=none] {x};
      \solonly{
      % Draw the angle
      \draw[thick,->] (0,0) -- (\rtheta{5}{0});
      \draw[thick,->] (0,0) -- (\rtheta{5}{300});

      % Draw angle arc
      \addplot[domain=0:300,-stealth, very thick] ({cos(x)},{sin(x)});

      % Label the angle
      \node at (\rtheta{1.5}{300/2}) {$\theta$};
      
       % Label the reference angle
      \node at (\rtheta{1.75}{340}) {\footnotesize$\pi/3$};
      }
    \end{axis}
\end{tikzpicture}&
\begin{tikzpicture}
    \begin{axis}[
        axis equal image,
        xmin=-5, xmax=5, ymin=-5, ymax=5,
        xlabel={$x$},
        ylabel={$y$},
        xtick=\empty, ytick=\empty,
        width=2.5in
      ]
      \addplot[draw=none] {x};
      \solonly{
      % Draw the angle
      \draw[thick,->] (0,0) -- (\rtheta{5}{0});
      \draw[thick,->] (0,0) -- (\rtheta{5}{225});

      % Draw angle arc
      \addplot[domain=0:225,-stealth, very thick] ({cos(x)},{sin(x)});

      % Label the angle
      \node at (\rtheta{1.5}{225/2}) {$\theta$};
      
      % Label the reference angle
      \node at (\rtheta{2}{200}) {\footnotesize$\pi/4$};
      }
    \end{axis}
\end{tikzpicture}\\

$\theta = \frac{5\pi}{6}$ \hspace{1em} &
$\theta = \frac{5\pi}{3}$ \hspace{1em} &
$\theta = \frac{5\pi}{4}$ \hspace{1em} \\[6pt]

reference angle $=$ \cfitb[\vphantom{\vspace{16pt}}$\frac{\pi}{6}$]{0.5in}{$\frac{\pi}{6}$} &
reference angle $=$ \cfitb[\vphantom{\vspace{16pt}}]{0.5in}{$\frac{\pi}{3}$} &
reference angle $=$ \cfitb[\vphantom{\vspace{16pt}}]{0.5in}{$\frac{\pi}{4}$}

\end{tabular}

\vspace{16pt}

\begin{tabular}{ccc}
(d)\hspace{1em} & (e)\hspace{1em} & (f)\hspace{1em} \\
\begin{tikzpicture}
    \begin{axis}[
        axis equal image,
        xmin=-5, xmax=5, ymin=-5, ymax=5,
        xlabel={$x$},
        ylabel={$y$},
        xtick=\empty, ytick=\empty,
        width=2.5in
      ]
      \addplot[draw=none] {x};
      \solonly{
      % Draw the angle
      \draw[thick,->] (0,0) -- (\rtheta{5}{0});
      \draw[thick,->] (0,0) -- (\rtheta{5}{210});

      % Draw angle arc
      \addplot[domain=210:360,stealth-, very thick] ({cos(x)},{sin(x)});

      % Label the angle
      \node at (\rtheta{1.5}{285}) {$\theta$};
      
      % Label the reference angle
      \node at (\rtheta{2}{195}) {\footnotesize$\pi/6$};
      }
    \end{axis}
\end{tikzpicture} &
\begin{tikzpicture}
    \begin{axis}[
        axis equal image,
        xmin=-5, xmax=5, ymin=-5, ymax=5,
        xlabel={$x$},
        ylabel={$y$},
        xtick=\empty, ytick=\empty,
        width=2.5in
      ]
      \addplot[draw=none] {x};
      \solonly{
      % Draw the angle
      \draw[thick,->] (0,0) -- (\rtheta{5}{0});
      \draw[thick,->] (0,0) -- (\rtheta{5}{60});

      % Draw angle arc
      \addplot[domain=60:360,stealth-, very thick] ({cos(x)},{sin(x)});

      % Label the angle
      \node at (\rtheta{1.5}{210}) {$\theta$};
      
      % Label the reference angle
      \node at (\rtheta{2}{25}) {\footnotesize$\pi/3$};
      }
    \end{axis}
\end{tikzpicture} &
\begin{tikzpicture}
    \begin{axis}[
        axis equal image,
        xmin=-5, xmax=5, ymin=-5, ymax=5,
        xlabel={$x$},
        ylabel={$y$},
        xtick=\empty, ytick=\empty,
        width=2.5in
      ]
      \addplot[draw=none] {x};
      \solonly{
      % Draw the angle
      \draw[thick,->] (0,0) -- (\rtheta{5}{0});
      \draw[thick,->] (0,0) -- (\rtheta{5}{450});

      % Draw angle arc
      \addplot[domain=0:450,-stealth, very thick] ({cos(x)*(0.5+x/900)},{sin(x)*(0.5+x/900)});

      % Label the angle
      \node at (\rtheta{1.5}{450/2}) {$\theta$};

      % Label the reference angle
      \node at (\rtheta{2}{45}) {\footnotesize$\pi/2$};
      }
    \end{axis}
\end{tikzpicture}\\

$\theta = -\frac{5\pi}{6}$ \hspace{1em} &
$\theta = -\frac{5\pi}{3}$ \hspace{1em} &
$\theta = \frac{5\pi}{2}$ \hspace{1em} \\[6pt]

reference angle $=$ \cfitb[\vphantom{\vspace{16pt}}]{0.5in}{$\frac{\pi}{6}$} &
reference angle $=$ \cfitb[\vphantom{\vspace{16pt}}]{0.5in}{$\frac{\pi}{3}$} &
reference angle $=$ \cfitb[\vphantom{\vspace{16pt}}]{0.5in}{$\frac{\pi}{2}$}

\end{tabular}
\end{center}

\item 
\begin{enumerate}
    \item 
\begin{problem}
Which quadrant does $\theta=\frac{100\pi}{3}$ land in? (It's efficient to use a reference angle.)
\end{problem}
\pspace{\vfill}

\begin{solution}
    The angle $\frac{100\pi}{3}$ is $\pi/3$ more than $\frac{99\pi}{3}=33\pi$. The terminal ray of $33\pi$ points in the negative $x$-direction, so an angle $\pi/3$ more than $33\pi$ must land in \fbox{quadrant III.}
\end{solution}

\item 
\begin{problem}
Which quadrant does $\theta=2$ (2 radians) land in?
\end{problem}
\pspace{\vfill}

\begin{solution}
    2 radians is less than $\pi$ radians but is more than $\pi/2$ radians, so the angle must land in \fbox{quadrant II.}
\end{solution}
    
\end{enumerate}
\pspace{\newpage}

\item 
\begin{problem}
Label the angles corresponding to the points shown below \textbf{in radians}. Use values between $0$ and $2\pi$.

\begin{center}
\begin{tabular}{cc}

\begin{tikzpicture}
      
  \begin{axis}[
  xmin=-1.1, xmax=1.1,
  ymin=-1.1, ymax=1.1,
  xlabel=\large$x$, ylabel=\large$y$,
  xtick=\empty, ytick=\empty,
  axis equal = true
  ]
  \coordinate (O) at (0,0);
  \coordinate (X) at (1,0);
  \coordinate (Y) at (0,1);
  \addplot[domain=0:360, data cs=polar] (x,1);
  \end{axis}
  \begin{scope}[x={($(X)-(O)$)}, y={($(Y)-(O)$)}, shift={(O)}]
    \foreach \A in {0,90,180,270} {
      \filldraw[radius=2pt] (\rtheta{1}{\A}) circle;
    }
    \foreach \A in {45, 135, 225, 315} {
      \filldraw[radius=2pt] (\rtheta{1}{\A}) circle;
    }
    \solonly{
    \node[above right] at (0:1) {$0$};
    \node[above right] at (45:1) {$\pi/4$};
    \node[above right] at (90:1) {$\pi/2$};
    \node[above left] at (135:1) {$3\pi/4$};
    \node[below left] at (180:1) {$\pi$};
    \node[below left] at (225:1) {$5\pi/4$};
    \node[below right] at (270:1) {$3\pi/2$};
    \node[below right] at (315:1) {$7\pi/4$};
    }
  \end{scope}
\end{tikzpicture}

&
\begin{tikzpicture}
      
  \begin{axis}[
  xmin=-1.1, xmax=1.1,
  ymin=-1.1, ymax=1.1,
  xlabel=\large$x$, ylabel=\large$y$,
  xtick=\empty, ytick=\empty,
  axis equal = true
  ]
  \coordinate (O) at (0,0);
  \coordinate (X) at (1,0);
  \coordinate (Y) at (0,1);
  \addplot[domain=0:360, data cs=polar] (x,1);
  \end{axis}
  \begin{scope}[x={($(X)-(O)$)}, y={($(Y)-(O)$)}, shift={(O)}]
    \foreach \A in {0,90,180,270} {
      \filldraw[radius=2pt] (\rtheta{1}{\A}) circle;
    }
    \foreach \A in {60, 150, 210, 330} {
      \filldraw[radius=2pt] (\rtheta{1}{\A}) circle;
    }
    \foreach \A in {30, 120, 240, 300} {
      \filldraw[radius=2pt] (\rtheta{1}{\A}) circle;
    }
    \solonly{
    \node[above right] at (0:1) {$0$};
    \node[above right] at (30:1) {$\pi/6$};
    \node[above right] at (60:1) {$\pi/3$};
    \node[above right] at (90:1) {$\pi/2$};
    \node[above left] at (120:1) {$2\pi/3$};
    \node[above left] at (150:1) {$5\pi/6$};
    \node[below left] at (180:1) {$\pi$};
    \node[below left] at (210:1) {$7\pi/6$};
    \node[below left] at (240:1) {$4\pi/3$};
    \node[below right] at (270:1) {$3\pi/2$};
    \node[below right] at (300:1) {$5\pi/3$};
    \node[below right] at (330:1) {$11\pi/6$};
    }
  \end{scope}
  
  \end{tikzpicture}
\end{tabular}
\end{center}
\end{problem}

% \item 
% \begin{problem}
%     On Problem Set 2, we investigated two types of \emph{special right triangles}. Find and label the missing side lengths of the special right triangles shown below, and label the angles in radians.
%          \begin{center}
%     \begin{tabular}{ccc}
%       an isosceles right triangle & & half of an equilateral triangle \\ \\
%       \begin{tikzpicture}[x=2.5cm, y=2.5cm]
%         \draw (0, 0) --  (1, 0) -- (1, 1) -- node[midway, above left] {1} (0,0);
%         \draw (1,0) rectangle +(-0.1, +0.1);
%       \end{tikzpicture}
%       &
%       \hspace{0.25in} &
%       \begin{tikzpicture}[x=1.5cm, y=1.5cm]
%         \draw (0, 0) -- (1, 0) -- (1, {sqrt(3)}) -- node[midway, above left] {1} (0,0);
%         \draw (1,0) rectangle +(-0.15, +0.15);
%       \end{tikzpicture} \\      
%     \end{tabular}  
%   \end{center}
% \end{problem}

\item 
\begin{problem}
The circle with radius 1 centered at the origin has a special name: the \textbf{unit circle.} Using what you know about special right triangles, label the coordinates of all the marked points on the unit circle. (Refer back to Problem Set 2 \#7 if needed.)
\begin{center}
\begin{tabular}{cc}

\begin{tikzpicture}
      
  \begin{axis}[
  xmin=-1.1, xmax=1.1,
  ymin=-1.1, ymax=1.1,
  xlabel=\large$x$, ylabel=\large$y$,
  xtick=\empty, ytick=\empty,
  axis equal = true
  ]
  \coordinate (O) at (0,0);
  \coordinate (X) at (1,0);
  \coordinate (Y) at (0,1);
  \addplot[domain=0:360, data cs=polar] (x,1);
    \draw (0,0) -- (0.707,0) -- (0.707,0.707) -- (0,0);
    \draw (0.707,0) rectangle (0.607, 0.1);
    \solonly{
    \draw[draw=none] (0,0) -- node[midway, below] {$\frac{\sqrt{2}}{2}$} (0.707,0) -- node[midway, right] {$\frac{\sqrt{2}}{2}$} (0.707,0.707) -- node[midway, above left] {$1$} (0,0);
    }
  \end{axis}
  \begin{scope}[x={($(X)-(O)$)}, y={($(Y)-(O)$)}, shift={(O)}]
    \foreach \A in {0,90,180,270} {
      \filldraw[radius=2pt] (\rtheta{1}{\A}) circle;
    }
    \foreach \A in {45, 135, 225, 315} {
      \filldraw[radius=2pt] (\rtheta{1}{\A}) circle;
    }
        \solonly{
    \node[above right] at (0:1) {$(1,0)$};
    \node[above right] at (45:1) {$\big(\frac{\sqrt{2}}{2},\frac{\sqrt{2}}{2}\big)$};
    \node[above right] at (90:1) {$(0,1)$};
    \node[above left] at (135:1) {$\big({-\frac{\sqrt{2}}{2}},\frac{\sqrt{2}}{2}\big)$};
    \node[below left] at (180:1) {$(-1,0)$};
    \node[below left] at (225:1) {$\big({-\frac{\sqrt{2}}{2}},{-\frac{\sqrt{2}}{2}}\big)$};
    \node[below right] at (270:1) {$(0,{-1})$};
    \node[below right] at (315:1) {$\big(\frac{\sqrt{2}}{2},{-\frac{\sqrt{2}}{2}}\big)$};
    }
  \end{scope}
  
\end{tikzpicture}

&
\begin{tikzpicture}
      
  \begin{axis}[
  xmin=-1.1, xmax=1.1,
  ymin=-1.1, ymax=1.1,
  xlabel=\large$x$, ylabel=\large$y$,
  xtick=\empty, ytick=\empty,
  axis equal = true
  ]
  \coordinate (O) at (0,0);
  \coordinate (X) at (1,0);
  \coordinate (Y) at (0,1);
  \addplot[domain=0:360, data cs=polar] (x,1);
  \draw (0,0) -- (0.5,0) -- (0.5,0.866) -- (0,0);
    \draw (0.5,0) rectangle (0.4, 0.1);
    \draw (0,0) -- (0.866,0) -- (0.866,0.5) -- (0,0);
  \draw (0.866,0) rectangle (0.766, 0.1);
  \solonly{
  \draw (0,0) -- node[midway, below] {$\frac{1}{2}$} (0.5,0) -- node[midway, right] {$\frac{\sqrt{3}}{2}$} (0.5,0.866) -- node[midway, above left] {$1$} (0,0);
  }
  \end{axis}
  \begin{scope}[x={($(X)-(O)$)}, y={($(Y)-(O)$)}, shift={(O)}]
    \foreach \A in {0,90,180,270} {
      \filldraw[radius=2pt] (\rtheta{1}{\A}) circle;
    }
    \foreach \A in {60, 150, 210, 330} {
      \filldraw[radius=2pt] (\rtheta{1}{\A}) circle;
    }
    \foreach \A in {30, 120, 240, 300} {
      \filldraw[radius=2pt] (\rtheta{1}{\A}) circle;
    }
\solonly{
    \node[above right] at (0:1) {$(1,0)$};
    \node[above right] at (30:1) {$\big(\frac{\sqrt{3}}{2},\frac{1}{2}\big)$};
    \node[above right] at (60:1) {$\big(\frac{1}{2},\frac{\sqrt{3}}{2}\big)$};
    \node[above right] at (90:1) {$(0,1)$};
    \node[above left] at (120:1) {$\big({-\frac{1}{2}},\frac{\sqrt{3}}{2}\big)$};
    \node[above left] at (150:1) {$\big({-\frac{\sqrt{3}}{2}},\frac{1}{2}\big)$};
    \node[below left] at (180:1) {$(-1,0)$};
    \node[below left] at (210:1) {$\big({-\frac{\sqrt{3}}{2}},{-\frac{1}{2}}\big)$};
    \node[below left] at (240:1) {$\big({-\frac{1}{2}},{-\frac{\sqrt{3}}{2}}\big)$};
    \node[below right] at (270:1) {$(0,{-1})$};
    \node[below right] at (300:1) {$\big({\frac{1}{2}},{-\frac{\sqrt{3}}{2}}\big)$};
    \node[below right] at (330:1) {$\big({\frac{\sqrt{3}}{2}},{-\frac{1}{2}}\big)$};
    }
  \end{scope}
  
  \end{tikzpicture}
\end{tabular}
\end{center}
\end{problem}


\end{enumerate}

\psreflection{
Optional reading for next class is \S 19.1 in Gottlieb.
}
\end{document}